\section{核心方法}
\label{sec:method}

本节阐述三项核心方法：分层决策、可插拔领域策略，以及三层嵌套自主循环。

\subsection{分层决策：规则优先、LLM 兜底}
\label{sec:method-decision}

决策的目标是：给定当前终端屏幕文本与上下文，输出一个结构化决策，说明代理是否需要外部动作、需要何种动作。WhatyTerm 采用两级决策架构。

\textbf{第一级：确定性规则状态机。} 一个纯规则函数按优先级从高到低对清洗后的屏幕文本做匹配，直接产出结构化决策，\emph{完全不调用 LLM}。其判定优先级大致为：运行中 $\to$ 评分/菜单界面 $\to$ 权限确认界面 $\to$ 致命错误 $\to$ shell 空闲提示符 $\to$ 代理空闲 $\to$ 其他。只有当规则无法定论（返回空）时，才进入第二级。

\textbf{第二级：LLM 兜底判官。} 对规则无法覆盖的开放场景，构造一份带优先级排序的判断规则提示词，强制模型只返回纯 JSON 决策。响应解析做了多重容错：剥离代码块标记、括号配对提取首个完整 JSON 对象、失败后重试修复，以应对模型输出不规范的情况。

无论哪一级，输出都遵循统一的动作模式（schema）：动作类型取自有限集合 $\{$文本输入, 菜单选择, 单字符, shell 命令, 确认, 警告, 自动修复, 无$\}$，下游据此用 \texttt{tmux send-keys} 发送对应按键。系统内置计数器分别统计“规则预判”与“LLM 分析”的次数，使两级占比可被量化（见 \S\ref{sec:eval}）。

\textbf{鲁棒状态识别。} 真实 TUI 屏幕是被 ANSI 转义序列、盒装 UI、滚动历史污染的脏文本，直接匹配极易误判。系统对 CSI/OSC/字符集切换等序列做多级清洗，并只对“最后若干行 / 最后若干字符”做窗口化匹配，避免历史缓冲区中的旧关键字干扰。一个典型的歧义消解案例是：Claude Code 空闲时底部仍常驻“esc to interrupt”状态栏，因此系统规定——只有在\emph{不存在}空闲提示符时该状态栏才计为“运行中”，而真正可靠的运行标志是进度计时指示器（形如 \texttt{...(Ns)}）。

\subsection{可插拔领域策略框架}
\label{sec:method-plugin}

不同任务类型对“何时自动继续、何时必须停手让人确认、何为错误”的判断规则截然不同。硬编码单一规则无法覆盖，为此系统提供基于策略模式的可插拔插件框架。每个插件实现统一接口：\texttt{matches}（用项目特征正则判断是否适配当前项目）、\texttt{detectPhase}（识别当前处于项目的哪个阶段）、\texttt{getPhaseConfig}（返回该阶段的自动动作、检查点、告警模式、提示词模板与是否需人工确认）、\texttt{analyzeStatus}（据阶段给出决策）。插件管理器按序注册十余个内置领域策略（全栈开发、论文写作、数据分析、TDD、安全审计等），通用兜底策略最后注册；选择时遍历非兜底插件返回首个匹配者，皆不中则用兜底。

该框架的两个设计要点值得强调。其一，\textbf{提示词与代码分离}：每个插件每个阶段的提示词以独立 Markdown 文件存放，支持片段复用与变量替换，并可热重载，构成一个模板化、可维护的领域知识库。其二，\textbf{阶段级自动化开关}：部分阶段（如“测试归档”）显式设置为\emph{禁止自动化、强制人工确认}，将“人在环边界”形式化为阶段配置的一部分，而非散落在代码中的特判。

\subsection{三层嵌套自主循环}
\label{sec:method-loop}

对于“完成一个完整需求”这类长程目标，单代理反复试错缺乏全局结构。WhatyTerm 将自主交付组织为三层嵌套循环，如表~\ref{tab:threelayer} 所示。

\begin{table}[t]
  \centering
  \small
  \caption{三层自主循环的职责与重试预算。}
  \label{tab:threelayer}
  \begin{tabular}{@{}lll@{}}
    \toprule
    层 & 职责 & 重试预算 \\
    \midrule
    macro & 需求$\to$PRD$\to$WBS$\to$接口契约；根因回退重拆 & 2 \\
    meso  & 子任务完成后集成验证（装依赖/编译/测试） & 3 \\
    micro & 单个子任务：执行$\to$验证$\to$失败重试 & 5 \\
    \bottomrule
  \end{tabular}
\end{table}

\textbf{macro 层：需求到可执行任务包。} 需求分解经三步流水线，每步一次 LLM 调用：澄清（判断需求是否充分，不足则产出澄清问题）、生成结构化 PRD、生成 WBS。WBS 将需求拆成子任务并为每个子任务附上\emph{接口契约}（输入、输出、验收标准、明确不包含的内容），要求“每任务规模受限、不跨层、依赖显式化、验收标准可机器验证”。随后按依赖关系对任务做拓扑排序，划分为可并行执行的批次，并处理循环依赖兜底。

\textbf{micro 层：Developer/Validator 双角色。} 每个子任务在 micro 层由两个职责分离的角色处理：Developer 只负责实现、自检、提交并输出可复用经验；Validator 只负责逐条对照验收标准判定、\emph{不修改代码}，并强制输出 \texttt{VALIDATION: PASS/FAIL} 及原因。这一“生成者/批判者分离”模式缓解了 LLM 自我裁判过宽的问题。子任务通过则标记完成，失败则计数重试，达上限则标记为 \texttt{blocked} 并跳过，防止卡死。算法~\ref{alg:micro} 给出 micro 循环。

\begin{algorithm}[t]
  \small
  \DontPrintSemicolon
  \KwIn{任务集合 $T$，每任务最大重试 $R{=}5$}
  \While{存在依赖已满足且未完成的任务 $t \gets \textsc{NextTask}(T)$}{
    $out \gets \textsc{RunDeveloper}(t)$\;
    $v \gets \textsc{RunValidator}(t, out)$\;
    \eIf{$v = \textsc{Pass}$}{
      标记 $t$ 为 completed；沉淀可复用经验\;
    }{
      $t.\mathit{retry} \gets t.\mathit{retry} + 1$\;
      \If{$t.\mathit{retry} \ge R$}{标记 $t$ 为 blocked（跳过）\;}
    }
  }
  \caption{micro 层子任务循环}
  \label{alg:micro}
\end{algorithm}

\textbf{meso/macro 层：集成验证与根因回退。} 所有子任务按批次并行执行完成后，meso 层自动探测并运行安装/编译/测试命令做集成验证。若验证失败，系统调用 LLM 对失败做\emph{语义分类}——归为接口不匹配、实现缺陷、设计缺陷或需求缺口之一——并据此选择\emph{最小修复范围}：接口/实现类问题在 meso 层局部重跑受影响的子任务；设计/需求类问题则回退到 macro 层重新拆解。这种“以失败类型选择修复范围”的策略，是本系统对通用代理循环的一个具体深化。三层各自的重试预算共同保证了整个流程的终止性。

\textbf{执行与完成检测。} 子任务以 headless 方式执行：将提示词写入临时文件，通过管道喂给 CLI 的非交互模式，并在命令末尾追加一个带返回码的完成标记；引擎轮询输出文件中的该标记与返回码来判定完成，而非从屏幕检测——这避免了命令回显中出现的标记字样造成误判。同时设置首字节、静默与总时长三重超时兜底。任务改动在专属 git 分支上进行以隔离影响，且支持每任务暂停等待人工审查。

\textbf{跨迭代经验沉淀。} Developer 在完成任务时可输出“可复用经验”，经进度管理器累积后注入后续任务的上下文，形成一个轻量的经验记忆机制，使后续任务复用已有约定、减少重复踩坑。所有进度以临时文件加原子重命名的方式持久化，支持断点续跑。
